% Question 2 content (input-only)
\textbf{Similarities.}
\begin{itemize}[leftmargin=*]
  \item \emph{Locality and parameter sharing}: Convolutional kernels apply local receptive fields with shared weights.
  \item \emph{Translation equivariance}: Feature maps shift consistently with input translations (exact under appropriate padding/stride settings).
\end{itemize}

\textbf{Differences.}
\begin{itemize}[leftmargin=*]
  \item \emph{Data domain}: 2D conv on H$\times$W (e.g., images); 3D conv on H$\times$W$\times$D/T (e.g., volumes/videos).
  \item \emph{Kernel sweep}: 2D slides over two axes; 3D slides over three axes, capturing volumetric or spatio-temporal patterns.
  \item \emph{Compute/params}: 3D kernels incur higher cost and memory; risk of overfitting increases under limited data.
  \item \emph{Use cases}: 2D for spatial cues (edges, textures); 3D for motion/shape continuity (temporal dynamics, 3D anatomy).
  \item \emph{Design variants}: Factorized (2+1)D or separable spatio-temporal convolutions reduce cost while preserving accuracy in videos.
\end{itemize}
